% !TeX root = ../algebra-1.tex
\chapter{Matrices and Linear Systems}\label{chap:algebra-matrices}

\section{Matrices, vectors and entrywise operations}
\begin{definition}
An $m\times n$ matrix over $F$ is a rectangular array with $m$ rows and
$n$ columns. Write $A=(a_{ij})$, where $a_{ij}$ is the entry in row $i$,
column $j$. A row vector has one row; a column vector has one column.
The set $F^n$ consists of column vectors of height $n$.
\end{definition}
For example,
\[
A=\begin{pmatrix}a&b&c\\d&e&f\end{pmatrix}
\]
is $2\times3$, its second row is $(d\ e\ f)$ and its third column is
$(c\ f)^{\mathsf T}$. The order of dimensions is always rows, then columns.
Two matrices are equal exactly when their sizes and corresponding entries agree.

Matrices of the same size add entry by entry, and scalar multiplication
multiplies every entry:
\[
(A+B)_{ij}=a_{ij}+b_{ij},\qquad(\lambda A)_{ij}=\lambda a_{ij}.
\]
For example,
\[
\begin{pmatrix}1&2\\4&5\end{pmatrix}
+\begin{pmatrix}0&1\\2&3\end{pmatrix}
=\begin{pmatrix}1&3\\6&8\end{pmatrix},\qquad
2\begin{pmatrix}1&-3\\0&1\end{pmatrix}
=\begin{pmatrix}2&-6\\0&2\end{pmatrix}.
\]
Addition is commutative and associative. For same-sized matrices,
\begin{align*}
(\alpha+\beta)A&=\alpha A+\beta A,&
\alpha(A+B)&=\alpha A+\alpha B,\\
\alpha(\beta A)&=(\alpha\beta)A.
\end{align*}
These identities follow entrywise from arithmetic in $F$.

\section{Matrix multiplication and linear combinations}
\begin{definition}[Matrix product]
If $A$ is $m\times n$ and $B$ is $n\times p$, their product is the
$m\times p$ matrix with
\[
(AB)_{ij}=\sum_{k=1}^n a_{ik}b_{kj}.
\]
The inner dimensions must agree. Each output entry is a row of $A$
multiplied into a column of $B$.
\end{definition}
\begin{example}[Sizes before arithmetic]
\[
\begin{pmatrix}a&b\\c&d\end{pmatrix}
\begin{pmatrix}e&f\\g&h\end{pmatrix}
=\begin{pmatrix}ae+bg&af+bh\\ce+dg&cf+dh\end{pmatrix}.
\]
A $2\times3$ matrix cannot multiply a $2\times1$ vector on its right.
But a $2\times3$ matrix can multiply a $3\times2$ matrix, with result
$2\times2$; for instance,
\[
\begin{pmatrix}1&2&3\\4&5&6\end{pmatrix}
\begin{pmatrix}1&2\\1&2\\1&2\end{pmatrix}
=\begin{pmatrix}6&12\\15&30\end{pmatrix}.
\]
\end{example}

\begin{definition}[Linear combination]
For same-sized matrices or vectors $v_1,\ldots,v_k$, an expression
$\sum_{j=1}^k\alpha_jv_j$, with coefficients $\alpha_j\in F$, is a linear
combination. Linear algebra studies which vectors can be produced this way
and whether their coefficients are uniquely determined.
\end{definition}

\begin{lemma}[Matrix times column; source Lemma 1.11]\label{lem:algebra-columns}
If $c_1,\ldots,c_n$ are the columns of an $m\times n$ matrix $A$ and
$x=(x_1,\ldots,x_n)^{\mathsf T}$, then
\[
Ax=x_1c_1+\cdots+x_nc_n.
\]
In particular, $Ae_j=c_j$, where $e_j$ is the standard vector with a single
$1$ in position $j$.
\end{lemma}
\begin{proof}
The $i$th entry of either side is $\sum_{j=1}^n a_{ij}x_j$.
For $x=e_j$, every coefficient except the $j$th is zero.
\end{proof}
\begin{example}
For $A=\begin{pmatrix}1&2\\3&4\end{pmatrix}$ and
$x=(1,-1)^{\mathsf T}$,
\[
Ax=\begin{pmatrix}1\\3\end{pmatrix}-\begin{pmatrix}2\\4\end{pmatrix}
=\begin{pmatrix}-1\\-1\end{pmatrix}.
\]
For the matrix with rows $(1,2,3),(4,5,6),(7,8,9)$, multiplication by
$e_2=(0,1,0)^{\mathsf T}$ selects its second column $(2,5,8)^{\mathsf T}$.
\end{example}

\begin{lemma}[Rowwise and columnwise multiplication]
If $B=[b_1\ \cdots\ b_p]$, then
$AB=[Ab_1\ \cdots\ Ab_p]$.
If $r_i$ is row $i$ of $A$, row $i$ of $AB$ is $r_iB$.
For a row $s=(s_1,\ldots,s_m)$ and rows $r_1,\ldots,r_m$ of $B$,
$sB=\sum_{i=1}^m s_ir_i$.
\end{lemma}
\begin{proof}
In each case, compare the $i,j$ entry using the product definition.
This supplies the rowwise/columnwise argument assigned as a problem in the source.
\end{proof}

\begin{example}[Multiplication encodes row and column changes]
Writing $A=\begin{pmatrix}a&b\\c&d\end{pmatrix}$,
\[
A\begin{pmatrix}0&1\\1&0\end{pmatrix}
=\begin{pmatrix}b&a\\d&c\end{pmatrix},\qquad
\begin{pmatrix}0&1\\1&1\end{pmatrix}A
=\begin{pmatrix}c&d\\a+c&b+d\end{pmatrix}.
\]
Right multiplication selects or combines columns; left multiplication
selects or combines rows.
\end{example}

\begin{proposition}[Algebra of products]
For compatible matrix sizes,
\begin{align*}
A(B+C)&=AB+AC,&(A+B)C&=AC+BC,\\
(\alpha A)B&=\alpha(AB)=A(\alpha B),&
A(BC)&=(AB)C.
\end{align*}
\end{proposition}
\begin{proof}
Distributivity and scalar compatibility follow entrywise. For associativity,
if $A$ is $m\times n$, $B$ is $n\times p$, and $C$ is $p\times q$,
\[
((AB)C)_{ij}=\sum_{\ell=1}^p\sum_{k=1}^n a_{ik}b_{k\ell}c_{\ell j}
=\sum_{k=1}^n\sum_{\ell=1}^p a_{ik}b_{k\ell}c_{\ell j}
=(A(BC))_{ij}.
\]
The finite double sums contain the same terms.
\end{proof}
Associativity allows unambiguous notation $ABC$. Commutativity generally fails:
\[
\begin{pmatrix}1&2\\3&4\end{pmatrix}
\begin{pmatrix}1&1\\1&1\end{pmatrix}
=\begin{pmatrix}3&3\\7&7\end{pmatrix},\qquad
\begin{pmatrix}1&1\\1&1\end{pmatrix}
\begin{pmatrix}1&2\\3&4\end{pmatrix}
=\begin{pmatrix}4&6\\4&6\end{pmatrix}.
\]

\section{Transpose, identity and zero matrices}
\begin{definition}[Transpose]
For $A=(a_{ij})$, its transpose $A^{\mathsf T}$ satisfies
$(A^{\mathsf T})_{ij}=a_{ji}$. It exchanges rows and columns and changes
size $m\times n$ to $n\times m$.
\end{definition}
\[
\begin{pmatrix}1&2&3\\4&5&6\end{pmatrix}^{\mathsf T}
=\begin{pmatrix}1&4\\2&5\\3&6\end{pmatrix}.
\]
The basic identities are
\begin{align*}
(A^{\mathsf T})^{\mathsf T}&=A,&
(A+B)^{\mathsf T}&=A^{\mathsf T}+B^{\mathsf T},\\
(\alpha A)^{\mathsf T}&=\alpha A^{\mathsf T},&
(AB)^{\mathsf T}&=B^{\mathsf T}A^{\mathsf T}.
\end{align*}
For the last identity, the $i,j$ entry of the left side is
$\sum_k a_{jk}b_{ki}$, equal to the $i,j$ entry of the right side.
The order reverses. Transpose alone does not conjugate complex entries.

\begin{definition}[Zero and identity]
The zero matrix $0_{m\times n}$ has every entry zero. The zero vector in
$F^n$ is $0_n$. The identity $I_n$ is $n\times n$ with $1$ on the diagonal
and $0$ elsewhere.
\end{definition}
For an $m\times n$ matrix $A$,
\[
A+0_{m\times n}=A,\qquad A I_n=A=I_mA.
\]
For square $A$, define $A^0=I_n$ and $A^k$ as the product of $k$ copies of
$A$ for $k\geq1$. Products involving zero matrices are zero with the
appropriate output size; similarly $A0_n=0_m$.

\section{Linear systems and augmented matrices}
A system of $m$ linear equations in $n$ variables has the form
\[
\sum_{j=1}^n a_{ij}x_j=b_i\quad(1\leq i\leq m).
\]
Put the coefficients into $A$, unknowns into $x$, and right sides into $b$.
The system is exactly $Ax=b$. For example,
\[
\begin{cases}2x+3y+4z=5,\\x+5z=0\end{cases}
\quad\Longleftrightarrow\quad
\begin{pmatrix}2&3&4\\1&0&5\end{pmatrix}
\begin{pmatrix}x\\y\\z\end{pmatrix}=\begin{pmatrix}5\\0\end{pmatrix}.
\]
The augmented matrix $[A\mid b]$ records coefficients and right sides together.
A solution is a vector $x$ satisfying every equation simultaneously.

\begin{definition}[Elementary row operations]
The allowed row operations are:
\begin{enumerate}
\item $R_j\leftarrow R_j+\lambda R_i$, where $i\ne j$;
\item $R_i\leftarrow\lambda R_i$, where $\lambda\ne0$;
\item $R_i\leftrightarrow R_j$.
\end{enumerate}
Their inverses subtract the same multiple, scale by $\lambda^{-1}$, and
swap the same rows, respectively.
\end{definition}
Apply operations to the \emph{entire} augmented matrix. Applying them only
to $A$ and leaving $b$ unchanged usually changes the system.

\begin{example}[Elimination with an augmented matrix]
For $3x+4y=6$, $x+2y=5$,
\begin{align*}
\left[\begin{array}{cc|c}3&4&6\\1&2&5\end{array}\right]
&\xrightarrow{R_1\leftarrow R_1-3R_2}
\left[\begin{array}{cc|c}0&-2&-9\\1&2&5\end{array}\right]\\
&\xrightarrow{R_1\leftarrow-\frac12R_1}
\left[\begin{array}{cc|c}0&1&9/2\\1&2&5\end{array}\right]\\
&\xrightarrow{R_2\leftarrow R_2-2R_1}
\left[\begin{array}{cc|c}0&1&9/2\\1&0&-4\end{array}\right]
\xrightarrow{R_1\leftrightarrow R_2}
\left[\begin{array}{cc|c}1&0&-4\\0&1&9/2\end{array}\right].
\end{align*}
Thus $x=-4$, $y=9/2$.
\end{example}

\section{Reduced row echelon form}
\begin{definition}[RREF]
The leading entry of a nonzero row is its leftmost nonzero entry.
A matrix is in reduced row echelon form (RREF) if:
\begin{enumerate}
\item every leading entry is $1$;
\item all zero rows are at the bottom;
\item a leading $1$ is the only nonzero entry in its column;
\item leading entries move strictly rightwards as one goes down nonzero rows.
\end{enumerate}
Columns containing leading entries are pivot columns.
\end{definition}
The zero matrix is in RREF. The matrix
$\begin{pmatrix}1&a&b&3\\0&0&1&-2\end{pmatrix}$ is in RREF precisely when
$b=0$. A row-echelon matrix need not have zeros \emph{above} its pivots;
RREF does.

\begin{theorem}[Existence and uniqueness of RREF]
Every matrix can be transformed by row operations into a unique RREF matrix.
\end{theorem}
\begin{proof}[Existence: elimination argument]
Scan columns from left to right. In the first available column with a nonzero
entry at or below the current pivot row, swap a nonzero entry into that row,
scale it to $1$, and eliminate all other entries in its column. Continue
below that pivot and to its right. New pivot rows have zeros in previous
pivot columns, so these operations preserve the earlier pivots.
Only finitely many rows and columns exist, and the final matrix satisfies
the RREF conditions. The source proves existence by an equivalent induction
on the number of columns. As in the source, uniqueness is stated without proof.
\end{proof}

\begin{example}[Source Example 1.34]\label{ex:algebra-rref}
Let
\[
A=\begin{pmatrix}1&2&3\\4&5&6\\7&8&9\end{pmatrix}.
\]
Subtract $4R_1$ from $R_2$ and $7R_1$ from $R_3$, then scale $R_2$ by
$-1/3$:
\[
A\sim\begin{pmatrix}1&2&3\\0&-3&-6\\0&-6&-12\end{pmatrix}
\sim\begin{pmatrix}1&2&3\\0&1&2\\0&-6&-12\end{pmatrix}
\sim\begin{pmatrix}1&0&-1\\0&1&2\\0&0&0\end{pmatrix}.
\]
The last step uses $R_1\leftarrow R_1-2R_2$ and $R_3\leftarrow R_3+6R_2$.
The pivot columns are $1$ and $2$. The symbol $\sim$ here means
row-equivalent, not equal.
\end{example}

\section{Reading the general solution}
In an augmented RREF matrix, a pivot in the last column represents an
impossible equation $0=1$. The system is then inconsistent.
Otherwise variables in coefficient pivot columns are \emph{pivot variables};
all remaining variables are \emph{free variables}.

\begin{theorem}[Free variables parameterise all solutions]
For a consistent RREF system, every assignment of scalars to the free
variables determines exactly one solution. Each nonzero row expresses its
pivot variable in terms of those free variables.
\end{theorem}
\begin{proof}
The pivot columns contain a single $1$ each. Hence each nonzero equation
has the form ``one pivot variable plus a linear combination of free variables
equals the right side''. No other equation contains that pivot variable.
Thus all pivots are uniquely determined, with no further restriction on the
free variables. Zero rows give only $0=0$.
\end{proof}

\begin{example}[Source Example 1.37]
For unknowns $(x,y,z,w)$, the augmented matrix
\[
\left[\begin{array}{cccc|c}
1&0&2&0&2\\0&1&4&0&3\\0&0&0&1&4\\0&0&0&0&0
\end{array}\right]
\]
gives $x+2z=2$, $y+4z=3$, $w=4$. Let $z=t$. Every solution is
\[
\begin{pmatrix}x\\y\\z\\w\end{pmatrix}
=\begin{pmatrix}2\\3\\0\\4\end{pmatrix}
+t\begin{pmatrix}-2\\-4\\1\\0\end{pmatrix},\qquad t\in F.
\]
Replacing the last augmented entry by $1$ produces a row $0=1$ and no solution.
\end{example}

\begin{example}[Source Example 1.39]
The system
\begin{align*}
2x_1+x_2+x_3-2x_4&=3,&x_1+x_3-2x_4&=2,\\
x_1+2x_3-4x_4&=1,&3x_1+x_2+3x_3-6x_4&=4
\end{align*}
reduces to $x_1=3$, $x_2=-2$, $x_3-2x_4=-1$. Its general solution is
\[
x=\begin{pmatrix}3\\-2\\-1\\0\end{pmatrix}
+t\begin{pmatrix}0\\0\\2\\1\end{pmatrix},\quad t\in F.
\]
Substitution into all four original equations is a useful check.
\end{example}

\begin{theorem}[More variables than equations; source Theorem 1.40]
If $A$ is $m\times n$ and $n>m$, then $Ax=0_m$ has a nonzero solution.
\end{theorem}
\begin{proof}
There are at most $m$ pivots, hence at least one free variable among the
$n$ unknowns. Give that free variable value $1$ and the other free variables
value $0$. The resulting solution is not the zero vector.
\end{proof}

\section{Inverses and elementary matrices}
\begin{definition}[Inverse]
An $n\times n$ matrix $A$ is invertible if a matrix $B$ satisfies
$AB=BA=I_n$. It is denoted $A^{-1}$ and is unique.
\end{definition}
Indeed, if $AB=I_n$ and $CA=I_n$, then
$B=(CA)B=C(AB)=C$. If $A$ is invertible, $Ax=b$ has the unique solution
$x=A^{-1}b$. The zero square matrix is not invertible, whereas $I_n$ is.
For compatible invertible square matrices,
\[
(AB)^{-1}=B^{-1}A^{-1},\qquad
(A^{\mathsf T})^{-1}=(A^{-1})^{\mathsf T}.
\]
These are checked by multiplication; inverse order reverses.

\begin{definition}[Elementary matrix]
An elementary matrix is obtained by applying one row operation to an identity
matrix. If the operation is $r$, its elementary matrix is $r(I_m)$.
\end{definition}
\begin{lemma}
For every $m\times n$ matrix $A$, performing $r$ is equivalent to left
multiplication by its elementary matrix: $r(A)=r(I_m)A$.
Elementary matrices are invertible, with inverses corresponding to the inverse
row operations. Consequently any row reduction has the form $R=EA$, where
$E$ is invertible.
\end{lemma}
\begin{proof}
Rows of $r(I_m)$ record exactly which linear combinations of the old rows
produce the new rows, so the rowwise multiplication lemma proves the identity.
Doing an inverse row operation undoes multiplication by $r(I_m)$.
For successive operations with elementary matrices $E_1,\ldots,E_k$,
$R=E_k\cdots E_1A$; the inverse product is $E_1^{-1}\cdots E_k^{-1}$.
\end{proof}
For example, $R_1\leftarrow R_1+2R_2$ for a three-row matrix is left
multiplication by
$\begin{pmatrix}1&2&0\\0&1&0\\0&0&1\end{pmatrix}$.

\begin{theorem}[Row operations preserve solutions]
If $[A'\mid b']=E[A\mid b]$ with $E$ invertible, then $Ax=b$ and
$A'x=b'$ have exactly the same solutions.
\end{theorem}
\begin{proof}
Here $A'=EA$ and $b'=Eb$. One equation follows from the other by
multiplying by $E$ or $E^{-1}$.
\end{proof}

\begin{theorem}[Little invertibility theorem; source Theorem 1.50]
For $A\in F^{n\times n}$, the following are equivalent:
$A$ is invertible; $Ax=0$ has only $x=0$; the RREF of $A$ is $I_n$.
\end{theorem}
\begin{proof}
An inverse makes $Ax=0$ imply $x=0$. If the kernel is zero, row reduction
can have no free variable, so each of the $n$ columns has a pivot and the
RREF is $I_n$. If $EA=I_n$ with $E$ invertible, then $A=E^{-1}$, so $A$
is invertible.
\end{proof}

\method{To compute an inverse, reduce $[A\mid I_n]$. If the left block can
be made $I_n$, the result is $[I_n\mid A^{-1}]$. If the left block has fewer
than $n$ pivots, $A$ is singular.}
The method works because $E[A\mid I_n]=[EA\mid E]$; when $EA=I_n$,
$E=A^{-1}$.

\begin{example}[Inverse by row reduction: source Example 1.52]
\begin{align*}
\left[\begin{array}{cc|cc}1&2&1&0\\3&4&0&1\end{array}\right]
&\sim\left[\begin{array}{cc|cc}1&2&1&0\\0&-2&-3&1\end{array}\right]\\
&\sim\left[\begin{array}{cc|cc}1&2&1&0\\0&1&3/2&-1/2\end{array}\right]
\sim\left[\begin{array}{cc|cc}1&0&-2&1\\0&1&3/2&-1/2\end{array}\right].
\end{align*}
Thus $\begin{pmatrix}1&2\\3&4\end{pmatrix}^{-1}
=\begin{pmatrix}-2&1\\3/2&-1/2\end{pmatrix}$.
For a general $2\times2$ matrix,
\[
\begin{pmatrix}a&b\\c&d\end{pmatrix}^{-1}
=\frac1{ad-bc}\begin{pmatrix}d&-b\\-c&a\end{pmatrix}
\quad\text{when }ad-bc\ne0.
\]
For larger matrices, elimination is usually more practical than an expanded
inverse formula. The determinant and cofactor formula appear in
Chapter~\ref{chap:algebra-eigenvalues}.
\end{example}
